\documentclass[12pt,a4paper]{article}
\usepackage[utf8]{inputenc}
\usepackage{graphicx}
\usepackage{color}
\usepackage[magyar]{babel}
\usepackage{listings}
\usepackage{float}

\usepackage{enumerate}
\usepackage{tikz}
\usetikzlibrary{shapes,arrows}
\usetikzlibrary{positioning}
\usetikzlibrary{arrows.meta}

\tikzset{
block/.style = {draw, fill=white, rectangle, minimum height=3em, minimum width=3em},
tmp/.style  = {coordinate}, 
input/.style = {coordinate},
output/.style= {coordinate},
pinstyle/.style = {pin edge={to-,thin,black}
}
}
\usepackage[siunitx]{circuitikz} 
\renewcommand{\lstlistingname}{Lista}
\usepackage[
	pdftitle={HA5KFU tanfolyam - Modulációs workshop},
	pdfauthor={HA5KFU Rádióamatőr Klub},
	pdfsubject={HA5KFU tanfolyam},
	pdfcreator={latex},
	pdfkeywords={ },
	pdfpagemode=UseOutlines,
	pdfdisplaydoctitle=true,
	pdflang=hu,
	unicode
]{hyperref}

\pagestyle{plain}
\sloppy
\begin{document}
\begin{center}
\includegraphics[width=300pt,keepaspectratio]{figures/ha5kfu.eps}
\\[0.5cm]
Rádióamatőr tanfolyamot segítő jegyzet, egyelőre kidolgozás alatt \\
Összeállította: Bazsó Márton, Szabó Áron % Feel free to add yourself
\\[1cm]

{\huge \bfseries Modulációs workshop \\[2cm]}



\end{center}

\renewcommand{\contentsname}{Tartalom}\tableofcontents 
\newpage

\newpage

\section{Bevezető}
A workshop célja bemutatni a különböző modulációs eljárások idő- és frekvenciatartománybeli viselkedését, illetve matematikai leírásukat. 

\section{Előkészületek}
A workshop Octave-ban és Matlab-ban is elvégezhető, a használt parancsok mindkettőben működőképesek.

\paragraph{Octave telepítés:} 
\begin{itemize}
	\item \textbf{Linux parancs:} \textit{sudo apt-get install octave}
	\item \textbf{Signal package telepítése:} Octave-on belül: \textit{pkg install -forge signal}\\
	Ha ez valamiért nem működne, akkor command line-ból is telepítheted: \textit{sudo apt-get install octave-signal}
\end{itemize}

\section{Feladatok}

\subsection{Szinusz jel generálása, ábrázolása}
Első feladatunk egy szinusz jel generálása lesz, melyet a \textit{sin()} függvénnyel tehetünk meg. Fontos, hogy ez a függvény egy adott érték szinuszát adja vissza, nekünk viszont egy időtartománybeli függvényre van szükségünk. Létre kell hoznunk tehát az "időt", amit jelöljünk \textit{t}-vel. Erre egy intervallumot jelölünk ki: \textit{t=0:0.05:10}. Itt 0 az intervallum kezdete, 10 a vége, amelyet 0.05-ös lépésközzel járunk végig (hasonlóképpen lehetne pl. \textit{t=0:10}, ekkor az alapértelmezett lépésköz 1).\\
Nézzük meg mit alkottunk a \textit{plot()} függvénnyel. Ha minden jól ment egy szinuszhoz hasonló jelalakot kapunk (ha nagyon torz, akkor lehet, hogy nem jól választottuk meg a lépésközt, vagy az intervallumot).

\end{document}
